\documentclass{article}
\usepackage{amsmath, amssymb, amsthm}

\setlength{\parindent}{0pt}

\newtheorem*{claim}{Claim}
\newtheorem{lemma}{Lemma}

\begin{document}

\begin{lemma}
    Let $X$ and $Y$ be subspaces of a vector space $V$. Then $X \cap Y$ is a subspace of $V$.
\end{lemma}

\begin{lemma}
    Let $V$ be a vector space. For $X, Y \subseteq V$ the sum $X + Y$ is the collection of all vectors
    $\mathbf{v}$ which can be represented as $\mathbf{v} = \mathbf{x} + \mathbf{y}$ for some 
    $\mathbf{x} \in X$ and $\mathbf{y} \in Y$. If $X$ and $Y$ are subspaces of $V$, then $X + Y$ is a subspace 
    of $V$.
\end{lemma}

\begin{claim}
    The smallest subspace of the space of $4 \times 4$ matrices that contains all upper triangular matrices 
    and all symmetric matrices is the set of all $4 \times 4$ matrices, and the largest subspace contained 
    in both of those subspaces is the set of all diagonal matrices.
\end{claim}

\begin{proof}
    Let $V$ be the space of all $4 \times 4$ matrices, and let $U$ be the subspace of all upper triangular 
    matrices and $S$ be the subspace of all symmetric matrices.
    
    \medskip

    We first show that the largest subspace contained in both $U$ and $S$ is $U \cap S$. Since $U$ and $S$ are 
    subspaces, by Lemma 1 $U \cap S$ is a subspace, and it is contained in both $U$ and $S$. Moreover, if $W$ is 
    any subspace with $W \subseteq U$ and $W \subseteq S$, then $W \subseteq U \cap S$ by definition of 
    intersection. Hence $U \cap S$ is the largest such subspace, and it remains to show that $U \cap S$ is the 
    set of all diagonal matrices. Let
    \[
    A = \begin{bmatrix}
        a_{11} & a_{12} & a_{13} & a_{14} \\
        a_{21} & a_{22} & a_{23} & a_{24} \\
        a_{31} & a_{32} & a_{33} & a_{34} \\
        a_{41} & a_{42} & a_{43} & a_{44}
    \end{bmatrix}
    \]
    where $A$ is an arbitrary matrix and $A \in U \cap S$. Upper triangularity annihilates everything below the 
    principal diagonal, giving 
    \[
    A = \begin{bmatrix}
        a_{11} & a_{12} & a_{13} & a_{14} \\
        0 & a_{22} & a_{23} & a_{24} \\
        0 & 0 & a_{33} & a_{34} \\
        0 & 0 & 0 & a_{44}
    \end{bmatrix}
    \]
    and symmetry requires that $a_{ij} = a_{ji}$, so for $i < j$ we get $a_{ij} = a_{ji} = 0$. Hence
    \[
    A = \begin{bmatrix}
        a_{11} & 0 & 0 & 0 \\
        0 & a_{22} & 0 & 0 \\
        0 & 0 & a_{33} & 0 \\
        0 & 0 & 0 & a_{44}
    \end{bmatrix}
    \]
    Thus every matrix in $U \cap S$ is diagonal. Conversely, let $D$ be an arbitrary diagonal matrix with 
    entries $d_{ij}$. All entries of $D$ below the principal diagonal are zero, so $D \in U$. Also 
    $d_{ij} = 0 = d_{ji}$ whenever $i \neq j$, so $D \in S$. Thus $D \in U \cap S$, and therefore $U \cap S$ is 
    exactly the set of all diagonal matrices.

    \medskip

    It remains to show that the smallest subspace containing both $U$ and $S$ is $V$. Let $X$ be any subspace 
    containing both $U$ and $S$. By closure under addition, $X$ contains $\mathbf{u} + \mathbf{s}$ for all 
    $\mathbf{u} \in U$ and $\mathbf{s} \in S$, so $U + S \subseteq X$. By Lemma 2, $U + S$ is a subspace of $V$, 
    and it contains both $U$ and $S$ since 
    $\mathbf{u} = \mathbf{u} + \mathbf{0}$ and $\mathbf{s} = \mathbf{0} + \mathbf{s}$. Hence $U + S$ is the 
    smallest subspace containing both $U$ and $S$, and it remains to show that $U + S = V$. Let
    \[
    B = \begin{bmatrix}
        a_{11} & a_{12} & a_{13} & a_{14} \\
        a_{21} & a_{22} & a_{23} & a_{24} \\
        a_{31} & a_{32} & a_{33} & a_{34} \\
        a_{41} & a_{42} & a_{43} & a_{44}
    \end{bmatrix}
    \]
    where $B$ is an arbitrary matrix. We must show that $B \in U + S$, i.e. $B = U_B + S_B$ for some matrices 
    $U_B \in U$ and $S_B \in S$. Let 
    \[
    S_B = \begin{bmatrix}
        0 & a_{21} & a_{31} & a_{41} \\
        a_{21} & 0 & a_{32} & a_{42} \\
        a_{31} & a_{32} & 0 & a_{43} \\
        a_{41} & a_{42} & a_{43} & 0
    \end{bmatrix}
    \]
    which is symmetric. Then 
    \[
    U_B = B - S_B = \begin{bmatrix}
        a_{11} & a_{12} - a_{21} & a_{13} - a_{31} & a_{14} - a_{41} \\
        0 & a_{22} & a_{23} - a_{32} & a_{24} - a_{42} \\
        0 & 0 & a_{33} & a_{34} - a_{43} \\
        0 & 0 & 0 & a_{44}
    \end{bmatrix}
    \]
    which is upper triangular. Therefore $B = (B - S_B) + S_B$ where $U_B = B - S_B \in U$ and $S_B \in S$, so 
    $B \in U + S$. Since $B$ was arbitrary, $V \subseteq U + S$, and combined with $U + S \subseteq V$ from 
    Lemma 2, we can conclude that $U + S = V$.

    \medskip

    Hence the smallest subspace containing both $U$ and $S$ is $V$ and the largest subspace contained in $U$ 
    and $S$ is the set of all diagonal matrices.
\end{proof}

\end{document}
